% Subsection 5.1.1: 实验环境与框架

本章在第四章固定的方法口径下执行 H1--H10 共 10 组评测。为保证结果只反映训练侧与推理侧的方法差异而不被运行环境污染，硬件、软件与并行粒度均在全部 10 组间保持一致。

\paragraph{（1）硬件与基座模型}
\begin{itemize}
    \item \textbf{GPU}：单张消费级/服务器级 GPU（训练与评测复用同一块显卡），显存不低于 24\,GB；实际运行中 \texttt{gradient\_checkpointing} 与 vLLM 的 \texttt{gpu\_memory\_utilization=0.9} 已在 7B 规模下联合限制峰值占用；
    \item \textbf{基座模型}：Qwen2.5-Coder-7B-Instruct，权重从官方仓库下载后直接落在 \texttt{models/Qwen2.5-Coder-7B-Instruct/}，10 组实验不换基座、不换权重 checkpoint；
    \item \textbf{数值精度}：训练与推理统一使用 \texttt{bf16}，避免因混合精度差异引入结果扰动。
\end{itemize}

\paragraph{（2）软件栈}
\begin{itemize}
    \item \textbf{训练侧}：\texttt{transformers}、\texttt{peft}（LoRA 实现）、\texttt{datasets}、\texttt{accelerate}，优化器固定为 \texttt{adamw\_torch}；
    \item \textbf{推理侧}：\texttt{vLLM}（PagedAttention + LoRA 子系统），负责批量生成与适配器挂载；
    \item \textbf{执行侧}：候选代码在独立 Python 子进程内调用 \texttt{scripts/validate\_code\_with\_test.py} 执行官方单测，\texttt{subprocess.run(timeout=5)} 负责超时隔离；
    \item \textbf{版本固定}：以仓库 \texttt{requirements.txt} 的 \texttt{pip freeze} 为准，详见 \texttt{SETUP.md}。
\end{itemize}

\paragraph{（3）训练-评测一键脚本}
全部 10 组实验通过 \texttt{scripts/train\_eval.sh} 串联：顶部的 \texttt{RUN\_TRAIN\_CODE / RUN\_TRAIN\_COT\_ZS / RUN\_TRAIN\_COT\_FS} 控制三类适配器的训练，\texttt{RUN\_H1 ~ RUN\_H10} 控制评测组别，\texttt{FS\_K} 同步约束 H3/H6/H9/H10 四组的示例数。该脚本保证「同一 \texttt{seed}、同一超参、同一评测口径」，使得复现者只需打开对应开关即可重跑任意子集。